% c13 原书 p.161–167（PDF p.168–174）
\begin{proposition}\label{it:2.2.7}
设 $K^{\cdot}$ 是带有递降滤过
\[
\cdots\supset F^pK^{\cdot}\supset F^{p+1}K^{\cdot}\supset\cdots
\]
的复形。若对每个 $n$，存在整数 $f(n)$ 与 $g(n)$，使得对一切 $p\le f(n)$ 有 $F^pK^n=K^n$，且对一切 $p\ge g(n)$ 有 $F^pK^n=0$，则上面构造的谱序列是双正则的。
\end{proposition}
\begin{proof}
给定 $p$ 与 $q$。若
\[
r>\max\bigl(g(p+q-1),g(p+q),g(p+q+1)\bigr)-p,
\]
则对 $n=p+q-1,p+q,p+q+1$ 有 $F^{p+r}K^n=0$，从而
\[
H^{p+q}(F^pK^{\cdot}/F^{p+r}K^{\cdot})\cong H^{p+q}(F^pK^{\cdot}).
\]
所以当 $r>\max\bigl(g(p+q-1),g(p+q),g(p+q+1)\bigr)-p$ 时，有 $Z_r^{pq}\cong Z_\infty^{pq}$。当
\[
r>p+1-\min\bigl(f(p+q-2),f(p+q-1),f(p+q)\bigr)
\]
时，对 $n=p+q-2,p+q-1,p+q$ 有 $F^{p-r+1}K^n=K^n$，从而
\[
H^{p+q-1}(F^{p-r+1}K^{\cdot}/F^pK^{\cdot})\cong H^{p+q-1}(K^{\cdot}/F^pK^{\cdot}).
\]
所以当 $r>p+1-\min\bigl(f(p+q-2),f(p+q-1),f(p+q)\bigr)$ 时，有 $B_r^{pq}\cong B_\infty^{pq}$。

对每个固定的 $n$，当 $p>g(n)$ 时有 $F^pK^n=0$，从而 $H^n(F^pK^{\cdot})=0$。所以当 $p>g(n)$ 时有 $F^pH^n=0$。当 $p<\min\bigl(f(n-1),f(n),f(n+1)\bigr)$ 时，对 $i=n-1,n,n+1$ 有 $F^pK^i=K^i$，从而 $H^n(F^pK^{\cdot})\cong H^n(K^{\cdot})$。所以当 $p<\min\bigl(f(n-1),f(n),f(n+1)\bigr)$ 时有 $F^pH^n=H^n$。
\end{proof}

设 $\cC$ 是阿贝尔范畴。$\cC$ 中的\emph{双复形}\idx{shuangfuxing}{双复形} $K^{\cdot\cdot}$ 是一族对象 $K^{ij}$（$i,j\in\bZ$）连同态射
\[
d_1^{ij}\colon K^{ij}\to K^{i+1,j},\qquad d_2^{ij}\colon K^{ij}\to K^{i,j+1},
\]
使得
\[
d_1^{i+1,j}d_1^{ij}=0,\quad d_2^{i,j+1}d_2^{ij}=0,\quad d_1^{i,j+1}d_2^{ij}=d_2^{i+1,j}d_1^{ij}.
\]
若 $\cC$ 中任意一族对象的直和都存在，或者对任意 $n$，满足 $i+j=n$ 且 $K^{ij}\ne 0$ 的对 $(i,j)$ 只有有限多个，则与 $K^{\cdot\cdot}$ 相伴的复形 $sK^{\cdot\cdot}$ 定义为
\[
(sK^{\cdot\cdot})^n=\bigoplus_{i+j=n}K^{ij},
\]
且 $d\colon(sK^{\cdot\cdot})^n\to(sK^{\cdot\cdot})^{n+1}$ 由
\[
dx=d_1x+(-1)^id_2x
\]
给出，其中 $x\in K^{ij}$。此时 $dd=0$。我们定义 $K^{\cdot\cdot}$ 的第 $n$ 个上同调对象 $H^n(K^{\cdot\cdot})$ 为 $sK^{\cdot\cdot}$ 的第 $n$ 个上同调对象。

对每个 $i$，
\[
\cdots\to K^{ij}\xrightarrow{d_2}K^{i,j+1}\to\cdots
\]
是复形。把它的第 $q$ 个上闭链（resp.\ 上边缘，resp.\ 上同调）对象记为 $Z_{II}^q(K^{i,\cdot})$（resp.\ $B_{II}^q(K^{i,\cdot})$，resp.\ $H_{II}^q(K^{i,\cdot})$）。注意
\[
\cdots\to Z_{II}^q(K^{i,\cdot})\xrightarrow{d_1}Z_{II}^q(K^{i+1,\cdot})\to\cdots,
\]
\[
\cdots\to B_{II}^q(K^{i,\cdot})\xrightarrow{d_1}B_{II}^q(K^{i+1,\cdot})\to\cdots,
\]
\[
\cdots\to H_{II}^q(K^{i,\cdot})\xrightarrow{d_1}H_{II}^q(K^{i+1,\cdot})\to\cdots
\]
都是复形。我们把复形 $H_{II}^q(K^{\cdot\cdot})$ 的第 $p$ 个上同调对象记为 $H_I^pH_{II}^q(K^{\cdot\cdot})$。类似地，有对象
\[
Z_I^p(K^{\cdot,j}),\qquad B_I^p(K^{\cdot,j}),\qquad H_I^p(K^{\cdot,j}),\qquad H_{II}^qH_I^p(K^{\cdot\cdot}).
\]

在 $sK^{\cdot\cdot}$ 上定义两个递降滤过如下：
\[
F_I^p(sK^{\cdot\cdot})^n=\bigoplus_{i+j=n,\ i\ge p}K^{ij},\qquad F_{II}^p(sK^{\cdot\cdot})^n=\bigoplus_{i+j=n,\ j\ge p}K^{ij}.
\]
这两个滤过定义了两个谱序列 $E_I$ 与 $E_{II}$。对于第一个滤过，有
\[
\bigl(F_I^p(sK^{\cdot\cdot})/F_I^{p+1}(sK^{\cdot\cdot})\bigr)^n\cong K^{p,n-p},
\]
\[
E_{I1}^{pq}=H^{p+q}\bigl(F_I^p(sK^{\cdot\cdot})/F_I^{p+1}(sK^{\cdot\cdot})\bigr)\cong H_{II}^q(K^{p,\cdot}).
\]
容易验证
\[
E_{I2}^{pq}=H_I^pH_{II}^q(K^{\cdot\cdot}).
\]
于是由第一个滤过定义的谱序列是
\[
E_{I2}^{pq}=H_I^pH_{II}^q(K^{\cdot\cdot})\Rightarrow H^{p+q}(K^{\cdot\cdot}).
\]
类似地，由第二个滤过定义的谱序列是
\[
E_{II2}^{pq}=H_{II}^pH_I^q(K^{\cdot\cdot})\Rightarrow H^{p+q}(K^{\cdot\cdot}).
\]
下面的命题由 \ref{it:2.2.7} 直接推出：

\begin{proposition}\label{it:2.2.8}
设 $K^{\cdot\cdot}$ 是双复形。设下列条件之一成立：
\begin{enumerate}[label=(\roman*),leftmargin=2.8em]
  \item 存在 $p_0,q_0$，使得只有当 $p\ge p_0$ 且 $q\ge q_0$ 时 $K^{pq}\ne 0$。
  \item 存在 $p_0,q_0$，使得只有当 $p\le p_0$ 且 $q\le q_0$ 时 $K^{pq}\ne 0$。
  \item 存在 $p_1\le p_2$，使得只有当 $p_1\le p\le p_2$ 时 $K^{pq}\ne 0$。
  \item 存在 $q_1\le q_2$，使得只有当 $q_1\le q\le q_2$ 时 $K^{pq}\ne 0$。
\end{enumerate}
则双复形 $K^{\cdot\cdot}$ 的两个谱序列都是双正则的。
\end{proposition}

\begin{proposition}\label{it:2.2.9}
设 $K^{\cdot}$ 是有足够内射对象的阿贝尔范畴中的复形。则存在双复形 $I^{\cdot\cdot}$，使得对每个 $q$，$I^{\cdot,q}$、$Z_{II}^q(I^{\cdot\cdot})$、$B_{II}^q(I^{\cdot\cdot})$ 与 $H_{II}^q(I^{\cdot\cdot})$ 分别是 $K^q$、$Z^q(K^{\cdot})$、$B^q(K^{\cdot})$ 与 $H^q(K^{\cdot})$ 的内射分解，且对每个 $i$，下列短正合序列都是分裂的：
\[
0\to B_{II}^q(I^{i,\cdot})\to Z_{II}^q(I^{i,\cdot})\to H_{II}^q(I^{i,\cdot})\to 0,
\]
\[
0\to Z_{II}^q(I^{i,\cdot})\to I^{i,q}\to B_{II}^{q+1}(I^{i,\cdot})\to 0.
\]
\end{proposition}

我们把 $I^{\cdot\cdot}$ 称为 $K^{\cdot}$ 的 \emph{Cartan-Eilenberg 分解}\idx{cartaneilenberg}{Cartan-Eilenberg 分解}。
\begin{tnote}
两条分裂短正合序列要合起来读：对每个 $i$，第一式给出该列关于 $q$ 的上同调分解，第二式给出该列本身由 $Z$ 与 $B$ 拼成的分解。随后的交换图中每一列（固定 $q$）自下而上的箭头表示 $K^q\to I^{0,q}\to I^{1,q}\to\cdots$ 是内射分解；横向箭头则表示这些内射分解之间由 $K^{\cdot}$ 的微分诱导的链映射。命题的要求正是：这些链映射同时把 $Z_{II}^q$、$B_{II}^q$、$H_{II}^q$ 也做成内射分解。[依据：原书 p.163 命题 2.2.9 及其所附交换图]
\end{tnote}
\[\begin{tikzcd}[row sep=1.5em,column sep=2.4em]
& \vdots & \vdots & \\
\cdots \arrow[r] & I^{1,q} \arrow[r] \arrow[u] & I^{1,q+1} \arrow[r] \arrow[u] & \cdots \\
\cdots \arrow[r] & I^{0,q} \arrow[r] \arrow[u] & I^{0,q+1} \arrow[r] \arrow[u] & \cdots \\
\cdots \arrow[r] & K^{q} \arrow[r] \arrow[u] & K^{q+1} \arrow[r] \arrow[u] & \cdots \\
& 0 \arrow[u] & 0 \arrow[u] &
\end{tikzcd}\]
\begin{proof}
对每个 $q$，取 $B^q(K^{\cdot})$ 的一个内射分解 $I_{B^q(K^{\cdot})}$ 与 $H^q(K^{\cdot})$ 的一个内射分解 $I_{H^q(K^{\cdot})}$。有正合序列
\[
0\to B^q(K^{\cdot})\to Z^q(K^{\cdot})\to H^q(K^{\cdot})\to 0.
\]
由 \ref{it:2.1.1}(iv)，存在内射分解 $I_{Z^q(K^{\cdot})}$ 以及复形的短正合序列
\[
0\to I_{B^q(K^{\cdot})}\to I_{Z^q(K^{\cdot})}\to I_{H^q(K^{\cdot})}\to 0,
\]
使得对每个 $i$ 有 $I_{Z^q(K^{\cdot})}^i=I_{B^q(K^{\cdot})}^i\oplus I_{H^q(K^{\cdot})}^i$。又有正合序列
\[
0\to Z^q(K^{\cdot})\to K^q\to B^{q+1}(K^{\cdot})\to 0.
\]
再次由 \ref{it:2.1.1}(iv)，存在 $K^q$ 的内射分解 $I^{\cdot,q}$ 以及复形的短正合序列
\[
0\to I_{Z^q(K^{\cdot})}\to I^{\cdot,q}\to I_{B^{q+1}(K^{\cdot})}\to 0,
\]
使得对每个 $i$ 有 $I^{i,q}=I_{Z^q(K^{\cdot})}^i\oplus I_{B^{q+1}(K^{\cdot})}^i$。于是双复形 $I^{\cdot\cdot}$ 具有所需的性质。
\end{proof}

\begin{proposition}\label{it:2.2.10}
设 $F\colon\cC\to\cC'$ 与 $G\colon\cC'\to\cC''$ 是阿贝尔范畴之间的两个左正合共变函子。设 $\cC$ 与 $\cC'$ 都有足够内射对象，且对 $\cC$ 中任意内射对象 $I$，对象 $F(I)$ 是 $G$-无环的。则对 $\cC$ 中每个对象 $A$，有双正则谱序列
\[
E_2^{pq}=\R^pG(\R^qF(A))\Rightarrow\R^{p+q}(GF)(A).
\]
\end{proposition}
\begin{proof}
设 $I^{\cdot}$ 是 $A$ 的内射分解，$J^{\cdot\cdot}$ 是 $F(I^{\cdot})$ 的 Cartan-Eilenberg 分解。考虑双复形 $G(J^{\cdot\cdot})$。对每个 $i$ 与每个 $q$，有分裂的短正合序列
\[
0\to B_{II}^q(J^{i,\cdot})\to Z_{II}^q(J^{i,\cdot})\to H_{II}^q(J^{i,\cdot})\to 0,
\]
\[
0\to Z_{II}^q(J^{i,\cdot})\to J^{i,q}\to B_{II}^{q+1}(J^{i,\cdot})\to 0.
\]
于是序列
\[
0\to G(B_{II}^q(J^{i,\cdot}))\to G(Z_{II}^q(J^{i,\cdot}))\to G(H_{II}^q(J^{i,\cdot}))\to 0,
\]
\[
0\to G(Z_{II}^q(J^{i,\cdot}))\to G(J^{i,q})\to G(B_{II}^{q+1}(J^{i,\cdot}))\to 0
\]
（分裂）正合。这说明
\[
H_{II}^q(G(J^{\cdot\cdot}))\cong G(H_{II}^q(J^{\cdot\cdot})).
\]
由于 $H_{II}^q(J^{\cdot\cdot})$ 是 $H^q(F(I^{\cdot}))\cong\R^qF(A)$ 的内射分解，有
\[
H_I^pH_{II}^q(G(J^{\cdot\cdot}))\cong H_I^p(G(H_{II}^q(J^{\cdot\cdot})))\cong\R^pG(\R^qF(A)).
\]
所以双复形 $G(J^{\cdot\cdot})$ 的第一个谱序列是
\[
E_{I2}^{pq}=\R^pG(\R^qF(A))\Rightarrow H^{p+q}(G(J^{\cdot\cdot})).
\]
由于 $J^{\cdot,p}$ 是 $F(I^p)$ 的内射分解，有
\[
H_I^q(G(J^{\cdot,p}))\cong\R^qG(F(I^p)).
\]
由于 $F(I^p)$ 是 $G$-无环的，有
\[
H_I^q(G(J^{\cdot,p}))\cong\R^qG(F(I^p))\cong
\begin{cases}
0, & \text{若 } q\ge 1,\\
GF(I^p), & \text{若 } q=0,
\end{cases}
\]
从而
\[
H_{II}^pH_I^q(G(J^{\cdot\cdot}))\cong
\begin{cases}
0, & \text{若 } q\ge 1,\\
H^p(GF(I^{\cdot}))\cong\R^p(GF)(A), & \text{若 } q=0.
\end{cases}
\]
所以双复形 $G(J^{\cdot\cdot})$ 的第二个谱序列
\[
E_{II2}^{pq}=H_{II}^pH_I^q(G(J^{\cdot\cdot}))\Rightarrow H^{p+q}(G(J^{\cdot\cdot}))
\]
退化。由 \ref{it:2.2.8}(i)，这两个谱序列都是双正则的。于是由 \ref{it:2.2.2}，对每个 $n$ 有
\[
\R^n(GF)(A)\cong H^n(G(J^{\cdot\cdot})).
\]
因此双复形 $G(J^{\cdot\cdot})$ 的第一个谱序列可以写成
\[
E_2^{pq}=\R^pG(\R^qF(A))\Rightarrow\R^{p+q}(GF)(A).
\]
\end{proof}

\begin{lemma}\label{it:2.2.11}
设 $f\colon X\to Y$ 是连续映射。对阿贝尔群的任意内射层 $\cI$，$f_*\cI$ 是内射的。
\end{lemma}
\begin{proof}
对 $Y$ 上阿贝尔群的任意层 $\cG$，有
\[
\Hom(\cG,f_*\cI)\cong\Hom(f^{-1}\cG,\cI).
\]
由于 $f^{-1}$ 是正合函子且 $\cI$ 是内射的，$\Hom(f^{-1}-,\cI)$ 是正合函子。所以 $\Hom(-,f_*\cI)$ 是正合函子，从而 $f_*\cI$ 是内射的。
\end{proof}

\begin{corollary}\label{it:2.2.12}
设 $f\colon X\to Y$ 是连续映射，$\cF$ 是 $X$ 上阿贝尔群的层。则有双正则谱序列
\[
E_2^{pq}=H^p(Y,\R^qf_*\cF)\Rightarrow H^{p+q}(X,\cF).
\]
\end{corollary}
\begin{proof}
把 \ref{it:2.2.10} 应用于函子 $F=f_*$ 与 $G=\Gamma(Y,-)$。我们有 $GF=\Gamma(X,-)$；由 \ref{it:2.2.11}，对 $X$ 上任意内射层 $\cI$，$F(\cI)$ 是 $G$-无环的。
\end{proof}

类似地有

\begin{corollary}\label{it:2.2.13}
设 $f\colon X\to Y$ 与 $g\colon Y\to Z$ 是连续映射，$\cF$ 是 $X$ 上阿贝尔群的层。则有双正则谱序列
\[
E_2^{pq}=\R^pg_*\R^qf_*\cF\Rightarrow\R^{p+q}(gf)_*\cF.
\]
\end{corollary}

\begin{lemma}\label{it:2.2.14}
设 $(X,\cO_X)$ 是环层空间，$\cI$ 是内射 $\cO_X$-模。则对任意 $\cO_X$-模 $\cF$，$\Hom_{\cO_X}(\cF,\cI)$ 是松弛的。
\end{lemma}
\begin{proof}
设 $U$ 是 $X$ 的开子集，$j\colon U\hookrightarrow X$ 是包含映射。给定态射 $\phi_U\colon\cF|_U\to\cI|_U$，我们要证明存在态射 $\phi\colon\cF\to\cI$ 使 $\phi|_U=\phi_U$。对 $U$ 上任意 $\cO_U$-模 $\cG$，回忆 $j_!\cG$ 是与下述预层相伴的 $X$ 上的层：
\[
V\mapsto
\begin{cases}
\cG(V), & \text{若 } V\subset U,\\
0, & \text{其他情形}.
\end{cases}
\]
有典范单射 $j_!(\cF|_U)\hookrightarrow\cF$。设 $\psi\colon j_!(\cF|_U)\to\cI$ 为复合
\[
j_!(\cF|_U)\xrightarrow{j_!(\phi_U)}j_!(\cI|_U)\to\cI.
\]
由于 $\cI$ 是内射的，态射 $\psi$ 可以延拓为态射 $\phi\colon\cF\to\cI$。此时 $\phi|_U=\phi_U$。
\end{proof}

\begin{proposition}\label{it:2.2.15}
设 $(X,\cO_X)$ 是环层空间，$\cF$ 与 $\cG$ 是 $\cO_X$-模。则有双正则谱序列
\[
E_2^{pq}=H^p(X,\Ext_{\cO_X}^q(\cF,\cG))\Rightarrow\Ext_{\cO_X}^{p+q}(\cF,\cG).
\]
\end{proposition}
\begin{proof}
把 \ref{it:2.2.10} 应用于函子
\[
F=\Hom_{\cO_X}(\cF,-),\qquad G=\Gamma(X,-).
\]
有 $GF=\Hom_{\cO_X}(\cF,-)$；由 \ref{it:2.2.14}，对任意内射 $\cO_X$-模 $\cI$，$F(\cI)$ 是 $G$-无环的。
\end{proof}

\begin{proposition}\label{it:2.2.16}
设 $(X,\cO_X)$ 是环层空间，$\cF$ 与 $\cG$ 是 $\cO_X$-模，$\cI$ 是 $\cG$ 的内射分解，且
\[
\cdots\to\cE_1\to\cE_0\to\cF\to 0
\]
是 $\cO_X$-模的正合序列，其中每个 $\cE_i$ 都是有限秩局部自由的。则有
\[
\Ext_{\cO_X}^n(\cF,\cG)\cong H^n(\Hom_{\cO_X}(\cF,\cI))\cong H^n(\Hom_{\cO_X}(\cE,\cG))\cong H^n(\Hom_{\cO_X}(\cE,\cI)).
\]
\end{proposition}
\begin{proof}
考虑双复形 $(K^{pq})=(\Hom_{\cO_X}(\cE_p,\cI^q))$。由于每个 $\cE_i$ 是有限秩局部自由的，每个函子 $\Hom_{\cO_X}(\cE_i,-)$ 都是正合的，从而
\[
H_{II}^q(K^{\cdot\cdot})=
\begin{cases}
0, & \text{若 } q\ge 1,\\
\Hom_{\cO_X}(\cE,\cG), & \text{若 } q=0.
\end{cases}
\]
于是
\[
H_I^pH_{II}^q(K^{\cdot\cdot})=
\begin{cases}
0, & \text{若 } q\ge 1,\\
H^p(\Hom_{\cO_X}(\cE,\cG)), & \text{若 } q=0.
\end{cases}
\]
所以双复形 $K^{\cdot\cdot}$ 的第一个谱序列退化。由 \ref{it:2.2.2}，这说明
\begin{equation*}\tag{1}
H^n(\Hom_{\cO_X}(\cE,\cG))\cong H^n(\Hom_{\cO_X}(\cE,\cI)).
\end{equation*}
由于每个 $\cI^q$ 是内射的，每个函子 $\Hom_{\cO_X}(-,\cI^q)$ 都是正合的，从而
\[
H_I^q(K^{\cdot\cdot})=
\begin{cases}
0, & \text{若 } q\ge 1,\\
\Hom_{\cO_X}(\cF,\cI), & \text{若 } q=0.
\end{cases}
\]
\begin{tnote}
此处原书把「对第一个指标取上同调」所得的对象记作 $H_I^q$，并把分类条件写成 $q\ge 1$ 与 $q=0$；而按 p.162 的记号 $H_I^pH_{II}^q$，第一个指标应写作 $p$。两种写法在此处指同一件事：固定 $q$ 后，沿 $p$ 方向的复形因 $\Hom_{\cO_X}(-,\cI^q)$ 正合而只在 $p=0$ 处有上同调 $\Hom_{\cO_X}(\cF,\cI^q)$。译文照原书印法照录，未改上标。[依据：原书 p.166；对照原书 p.162 的 $H_I^pH_{II}^q$ 定义]
\end{tnote}
于是
\[
H_{II}^pH_I^q(K^{\cdot\cdot})=
\begin{cases}
0, & \text{若 } q\ge 1,\\
\Ext_{\cO_X}^p(\cF,\cG), & \text{若 } q=0.
\end{cases}
\]
所以双复形 $K^{\cdot\cdot}$ 的第二个谱序列退化。由 \ref{it:2.2.2}，这说明
\begin{equation*}\tag{2}
\Ext_{\cO_X}^n(\cF,\cG)\cong H^n(\Hom_{\cO_X}(\cE,\cI)).
\end{equation*}
结合 (1) 与 (2)，得到
\[
\Ext_{\cO_X}^n(\cF,\cG)\cong H^n(\Hom_{\cO_X}(\cE,\cI))\cong H^n(\Hom_{\cO_X}(\cE,\cG)).
\]
\end{proof}

下面命题的证明留给读者：

\begin{proposition}\label{it:2.2.17}
设 $A$ 是环，$M$ 与 $N$ 是 $A$-模，且
\[
\cdots\to E^{-1}\to E^0\to M\to 0,\qquad\cdots\to F^{-1}\to F^0\to N\to 0
\]
是正合序列，其中 $E^i$ 与 $F^i$ 都是平坦 $A$-模。则有
\[
\Tor_n^A(M,N)\cong H^{-n}(M\tensor_A F^{\cdot})\cong H^{-n}(E^{\cdot}\tensor_A N)\cong H^{-n}(E^{\cdot}\tensor_A F^{\cdot}).
\]
特别地，有 $\Tor_n^A(M,N)\cong\Tor_n^A(N,M)$。
\end{proposition}

\begin{exercises}
  \item \label{ex:2.2.1}
  设 $A$ 是交换环。对任意 $A$-模 $M$，令 $N\mapsto\Ext_A^i(M,N)$ 为 $A$-模范畴上函子 $N\mapsto\Hom_A(M,N)$ 的右导出函子。对 $\Spec A$ 的任意仿射开子集 $U$，令
  \[
  A_U=\Gamma(U,\cO_{\Spec A}),\qquad M_U=\Gamma(U,M^{\sim}),\qquad N_U=\Gamma(U,N^{\sim}).
  \]
  设 $A$ 是诺特的且 $M$ 是有限生成的。证明
  \[
  \Ext_{A_U}^i(M_U,N_U)\cong\Ext_{\cO_U}^i(M^{\sim}|_U,N^{\sim}|_U).
  \]
  \item \label{ex:2.2.2}
  设 $R$ 是主理想整环。
  \begin{enumerate}[label=(\arabic*),leftmargin=2.2em]
    \item 对每个有限生成 $R$-模 $M$、每个 $R$-模 $N$ 与每个 $i\ne 0,1$，有
    \[
    \Ext_R^i(M,N)=0,\qquad \Tor_i^R(M,N)=0.
    \]
    \item 设 $K^{\cdot}$ 是有限秩自由 $R$-模的复形。证明有短正合序列
    \[
    0\to H^n(K^{\cdot})\tensor_R N\to H^n(K^{\cdot}\tensor_R N)\to\Tor_1^R(H^{n+1}(K^{\cdot}),N)\to 0,
    \]
    \[
    0\to\Ext_R^1(H^{-(n-1)}(K^{\cdot}),N)\to H^n(\Hom_R(K,N))\to\Hom_R(H^{-n}(K^{\cdot}),N)\to 0.
    \]
  \end{enumerate}
  \item \label{ex:2.2.3}
  设 $X$ 是概形，$\cF$ 是有限秩局部自由的 $\cO_X$-模，$\cG$ 是 $\cO_X$-模。证明对一切 $i$ 有同构
  \[
  H^i(X,\Hom_{\cO_X}(\cF,\cG))\cong\Ext_{\cO_X}^i(\cF,\cG).
  \]
  \item \label{ex:2.2.4}
  本题需要代数拓扑的知识。设 $X$ 是拓扑空间。回忆：奇异同调 $H_i(X)$ 定义为奇异链复形 $S.(X)$ 的第 $i$ 个同调群。对 $X$ 的任意子空间 $Y$，相对同调 $H_i(X,Y)$ 定义为复形 $S.(X)/S.(Y)$ 的第 $i$ 个同调群。
  \begin{enumerate}[label=(\arabic*),leftmargin=2.2em]
    \item 对任意三元组 $Z\subset Y\subset X$，有相对同调群的长正合序列
    \[
    \cdots\to H_i(Y,Z)\to H_i(X,Z)\to H_i(X,Y)\xrightarrow{\partial}H_{i-1}(Y,Z)\to\cdots.
    \]
    \item 设有一列子空间
    \[
    X_0\subset X_1\subset\cdots
    \]
    使得 $X=\bigcup_p X_p$。取 $X_i=\emptyset$（$i<0$）。设 $X$ 的每个紧子集都包含在某个 $X_p$ 中。构造一个正则谱序列
    \[
    E_1^{pq}\Rightarrow H^n,
    \]
    使得 $H^n=H_{-n}(X)$ 且 $E_1^{pq}=H_{p-q}(X_{-p},X_{-p-1})$。
    \item 进而假设当 $i\ne p$ 时 $H_i(X_p,X_{p-1})$ 为零。证明 $H_i(X)$ 同构于复形
    \[
    \cdots\to H_p(X_p,X_{p-1})\xrightarrow{\partial}H_{p-1}(X_{p-1},X_{p-2})\to\cdots
    \]
    的第 $i$ 个同调群。
  \end{enumerate}
\end{exercises}

